% Chapter 2: Data Preprocessing
\chapter{Data Preprocessing}

\section{Overview}
This chapter describes the preprocessing pipeline applied to each dataset before integration. The preprocessing decisions were guided by data quality issues identified during EDA (Chapter 1) and requirements of the machine learning algorithms.

\section{Fire Dataset Preprocessing}

The fire dataset from MODIS/VIIRS satellites required several cleaning steps to ensure data quality.

\subsection{Fire Type Filtering}
Only confirmed vegetation fires (type 0) were retained; presumed fires (type 2) were excluded to reduce false positives. This filtering removed ambiguous detections that could introduce noise.

\subsection{Data Cleaning}
\begin{itemize}
    \item \textbf{Duplicates:} Removed duplicate records (identical coordinates and timestamps)
    \item \textbf{Missing values:} Numeric columns imputed with median; categorical columns with mode
    \item \textbf{Latitude constraint:} Points below 33°N excluded (outside Algeria/Tunisia boundaries)
    \item \textbf{Column selection:} Only \texttt{longitude}, \texttt{latitude}, and \texttt{class} retained
\end{itemize}

\subsection{Outlier Treatment}
Outliers detected using the IQR method (1.5$\times$IQR beyond Q1/Q3). Outlier values were replaced with the column median to preserve data volume while reducing extreme value influence.

\subsection{Negative Sample Generation}
To create a balanced classification dataset, artificial non-fire points (class 0) were generated:
\begin{itemize}
    \item Random coordinates within the fire points' bounding box
    \item 1.5$\times$ the number of fire points generated
    \item Latitude constraint ($\geq$33°N) respected
    \item Combined dataset shuffled with fixed random seed for reproducibility
\end{itemize}

\section{Climate Dataset Preprocessing}

WorldClim data (precipitation, maximum and minimum temperature) was processed as monthly rasters.

\subsection{Raster Processing}
\begin{itemize}
    \item \textbf{Clipping:} Rasters clipped to Algeria and Tunisia boundaries using dissolved country shapefiles
    \item \textbf{Temporal alignment:} Files matched by date pattern (\texttt{YYYY-MM}) to ensure consistent temporal coverage
    \item \textbf{Unit conversion:} Temperature values scaled (divided by 10) to Celsius; precipitation in mm/month
\end{itemize}

\subsection{Feature Extraction at Fire Points}
For each fire location, climate values were extracted using raster coordinate transformation (\texttt{rowcol}). Seasonal aggregations computed:
\begin{itemize}
    \item Winter (DJF), Spring (MAM), Summer (JJA), Autumn (SON)
    \item Mean precipitation and temperature per season
\end{itemize}

\subsection{Normalization}
Z-score normalization applied to all climate variables:
\begin{equation}
    z = \frac{x - \mu}{\sigma}
\end{equation}

\section{Elevation Dataset Preprocessing}

The GMTED2010 DEM (15 arc-seconds, $\sim$450m resolution) was processed to extract topographic features.

\subsection{Data Quality}
NoData values ($-32768$) were masked during processing. The raster was clipped to a buffered bounding box around fire points for efficiency.

\subsection{Derived Terrain Features}
Four terrain variables were computed from the elevation raster:

\begin{itemize}
    \item \textbf{Slope:} Gradient magnitude in degrees, computed from x/y partial derivatives:
    \begin{equation}
        \text{slope} = \arctan\left(\sqrt{\left(\frac{\partial z}{\partial x}\right)^2 + \left(\frac{\partial z}{\partial y}\right)^2}\right) \times \frac{180}{\pi}
    \end{equation}
    
    \item \textbf{Aspect:} Slope orientation (0--360°, compass bearing with 0°=North)
    
    \item \textbf{Terrain Roughness Index (TRI):} Root mean square deviation from neighboring cells in a 3$\times$3 window---measures local terrain variability
    
    \item \textbf{Roughness:} Alternative measure using local standard deviation
\end{itemize}

\subsection{Outlier Treatment and Normalization}
IQR-based outlier replacement applied to elevation features, followed by StandardScaler normalization.

\section{Soil Dataset Preprocessing}

The HWSD2 database consists of an MDB file (tabular attributes) and a raster file (spatial mapping).

\subsection{Data Loading and Spatial Operations}
\begin{itemize}
    \item Loaded \texttt{HWSD2\_LAYERS} table, filtered for D1 layer (topsoil, 0--30cm)
    \item 22 features selected: 8 physical (COARSE, SAND, SILT, CLAY, textures, BULK) and 14 chemical (ORG\_CARBON, PH\_WATER, CEC, etc.)
    \item Raster clipped to fire bounding box; soil unit IDs extracted at fire coordinates
\end{itemize}

\subsection{Sentinel Value Handling}
HWSD2 uses $-9$ and similar values as missing data indicators. These sentinels were replaced with NaN before imputation.

\subsection{Missing Value Imputation}
\begin{itemize}
    \item \textbf{Numeric columns:} Median imputation
    \item \textbf{Categorical columns:} Mode imputation
    \item \textbf{Duplicate handling:} All rows retained (same soil properties can map to multiple geographic units)
\end{itemize}

\subsection{Normalization}
Z-score normalization (StandardScaler) applied to all 20 numeric soil features. Categorical texture variables (\texttt{TEXTURE\_USDA}, \texttt{TEXTURE\_SOTER}) processed separately via one-hot encoding.

\subsection{One-Hot Encoding}
The \texttt{TEXTURE\_SOTER} categorical variable was one-hot encoded to create binary columns for each texture class, enabling use in numerical algorithms.

\section{Data Integration}

\subsection{Spatial Alignment}
All datasets were aligned to WGS84 (EPSG:4326). Fire points served as the spatial reference for extracting values from raster sources.

\subsection{Merging Strategy}
The integration followed a fire-centric approach with left joins:
\begin{enumerate}
    \item \texttt{fires.csv} as base dataset (longitude, latitude, class)
    \item Merge climate features on (longitude, latitude, class)
    \item Merge elevation features on (longitude, latitude, class)
    \item Merge soil features on (longitude, latitude, class)
\end{enumerate}

\begin{figure}[H]
    \centering
    \fbox{\parbox{0.85\textwidth}{\centering\vspace{3cm}\textbf{[PLACEHOLDER: Data Integration Flowchart]}\vspace{3cm}}}
    \caption{Data integration and merging workflow}
    \label{fig:merging_flowchart}
\end{figure}

\section{Final Dataset}

The merged dataset combines all preprocessed features:

\begin{table}[H]
\centering
\caption{Final Merged Dataset Structure}
\label{tab:merged_structure}
\begin{tabular}{|l|l|c|}
\hline
\textbf{Category} & \textbf{Variables} & \textbf{Count} \\
\hline
Location & longitude, latitude & 2 \\
Climate & precipitation, tmax, tmin (seasonal aggregates) & 12+ \\
Topography & elevation, slope, aspect, TRI, roughness & 5 \\
Soil & physical (SAND, SILT, CLAY, BULK...) + chemical (ORG\_CARBON, PH...) & 22 \\
Texture & one-hot encoded TEXTURE\_SOTER categories & 8+ \\
Target & class (binary: 0=no fire, 1=fire) & 1 \\
\hline
\textbf{Total} & & \textbf{50+} \\
\hline
\end{tabular}
\end{table}

\subsection{Class Balance}
The final dataset maintains a 1:1.5 ratio (fire:non-fire), providing sufficient negative samples while avoiding extreme imbalance that could bias classifiers toward the majority class.
